\chapter{主要研究内容及研究方案}
根据上述分析可知，目前关于non-IID数据和标签噪声对SplitFed影响的研究较为稀缺，本课题可以填补这一空白。本课题将重点关注下面两个问题：
1.SplitFed在低质量数据场景下的收敛证明。
2.提出基于SplitFed的对低质量数据具有鲁棒性的分布式算法。
\section{收敛性分析}
\subsection{Non-IID 数据}
目前对于non-IID数据下分布式框架的收敛性分析主要集中在联邦学习上，SplitFed与联邦学习的主要区别是联邦学习的模型在不同用户间进行完整的训练，而SplitFed的模型的不同部分会分别在用户和计算服务器上进行训练。因此，建立SplitFed中被拆分的模型与完整模型之间的联系将是SplitFed收敛性证明的关键。
\subsection{标签噪声}
目前有成熟的理论来衡量分布式学习中用户数据的分布和质量，但这些标准没有考虑标签噪声。因此需要探索标签噪声如何影响用户数据的质量和分布，并使用目前的主流理论衡量标签噪声的影响，进而证明噪声标签如何影响SplitFed的收敛性。

\section{基于SplitFed的算法研究}
\subsection{non-IID数据}
该部分研究旨在探索具有鲁棒性的SplitFed算法，针对non-IID数据进行优化。通过分析non-IID数据对SplitFed的影响，设计改进算法以提高其在这种数据环境下的收敛性和性能。通过引入新的损失函数、优化聚合策略以及调整本地更新步骤，增强SplitFed在低质量和非一致性数据上的鲁棒性。
\subsection{标签噪声}
该部分研究旨在探索对包含标签噪声数据具有鲁棒性的SplitFed算法。通过分析标签噪声对SplitFed性能的影响，识别在模型训练中可能出现误导性更新的环节。随后，通过引入噪声建模、鲁棒性损失函数和自蒸馏技术来设计对标签噪声具有鲁棒性的SplitFed算法，提高模型对标签噪声的适应能力。此外，不同类型和程度的标签噪声都会被纳入考虑，对设计出来的算法进行系统性实验，评估改进算法的效果与稳定性。